% T2 candidate (sonnet3) for fig:staging — same gallery-filling cartoon as v7-11 (structural
% information preserved), now tied to a true-to-scale voltage axis carrying the four staging
% transitions' initial U values and a REAL dQ/dV peak glyph at each (eq:belliden bell,
% height Q_j/(4w_j) from eq:eqpeak, values transcribed from table tab:staging and the width
% fallback list — all coordinates from T2_calc.py, not freehand).
\begin{tikzpicture}[x=1.05cm,y=0.5cm,
  lab/.style={font=\scriptsize,align=center},
  fill4/.style={fill=black!72},fillL/.style={fill=black!24}]
% ---- (top) five stage columns: graphene layers + Li-filled galleries, unchanged structure ----
\def\colw{0.85}
\foreach \x/\name/\filled/\style in {0/{stage 4}/{1}/fill4, 1.4/{stage 3}/{1}/fill4, 2.8/{stage 2L}/{all}/fillL, 4.2/{stage 2}/{all}/fill4, 5.6/{stage 1}/{dense}/fill4}{
  \draw[semithick] (\x,0)--(\x+\colw,0);
  \foreach \k in {1,2,3,4,5,6}{\draw[semithick] (\x,0.6*\k)--(\x+\colw,0.6*\k);}
  \node[lab,below] at (\x+\colw/2,-0.25) {\name};
}
\foreach \k in {1,5}{\fill[black!72] (0,0.6*\k-0.42) rectangle (0+\colw,0.6*\k-0.18);}
\foreach \k in {1,4}{\fill[black!72] (1.4,0.6*\k-0.42) rectangle (1.4+\colw,0.6*\k-0.18);}
\foreach \k in {1,3,5}{\fill[black!24] (2.8,0.6*\k-0.42) rectangle (2.8+\colw,0.6*\k-0.18);}
\foreach \k in {1,3,5}{\fill[black!72] (4.2,0.6*\k-0.42) rectangle (4.2+\colw,0.6*\k-0.18);}
\foreach \k in {1,2,3,4,5,6}{\fill[black!72] (5.6,0.6*\k-0.42) rectangle (5.6+\colw,0.6*\k-0.18);}
\draw[-{Latex[length=1.8mm]},thick] (0,4.05)--(6.45,4.05) node[pos=0.5,above,font=\scriptsize] {lithiation (charge): $4\to3\to2\mathrm L\to2\to1$};
% ---- schematic transition boundaries (even column pitch) dropped down to the true-scale axis ----
\draw[densely dotted] (1.125,-0.5) -- (1.015,-2.3);
\draw[densely dotted] (2.525,-0.5) -- (3.518,-2.3);
\draw[densely dotted] (3.925,-0.5) -- (4.233,-2.3);
\draw[densely dotted] (5.325,-0.5) -- (5.485,-2.3);
\node[font=\tiny\itshape,align=center] at (3.2,-0.85) {schematic column pitch above $\ne$ true voltage spacing below};
% ==== (bottom) true-to-scale voltage axis with real dQ/dV peak glyphs ====
\begin{scope}[yshift=-1.15cm]
\draw[-{Latex[length=1.6mm]},thick] (-1.75,0) -- (7.35,0) node[below,font=\scriptsize] {$U$ [V] (staging transition, initial value; table)};
\foreach \x/\lab in {1.015/0.210, 3.518/0.140, 4.233/0.120, 5.485/0.085}{
  \draw (\x,0.03) -- (\x,-0.03);
  \node[font=\tiny,below] at (\x,-0.06) {\lab};
}
% dQ/dV bells: eq:belliden shape, height Q_j/(4w_j); fill opacity so overlaps add visually
% (matches eq:sum — the real curve is the sum of these four components)
\fill[black!22,fill opacity=0.55] (-1.416,0) -- plot[smooth] coordinates {
(-1.416,0.0240) (-1.273,0.0289) (-1.130,0.0347) (-0.987,0.0415) (-0.844,0.0494) (-0.701,0.0586) (-0.558,0.0690) (-0.415,0.0806) (-0.272,0.0935) (-0.129,0.1073) (0.014,0.1219) (0.157,0.1366) (0.300,0.1510) (0.443,0.1643) (0.586,0.1757) (0.729,0.1845) (0.872,0.1901) (1.015,0.1920) (1.158,0.1901) (1.301,0.1845) (1.444,0.1757) (1.587,0.1643) (1.730,0.1510) (1.873,0.1366) (2.016,0.1219) (2.159,0.1073) (2.302,0.0935) (2.445,0.0806) (2.588,0.0690) (2.732,0.0586) (2.875,0.0494) (3.018,0.0415) (3.161,0.0347) (3.304,0.0289) (3.447,0.0240)
} -- (3.447,0) -- cycle;
\draw[semithick] plot[smooth] coordinates {
(-1.416,0.0240) (-1.273,0.0289) (-1.130,0.0347) (-0.987,0.0415) (-0.844,0.0494) (-0.701,0.0586) (-0.558,0.0690) (-0.415,0.0806) (-0.272,0.0935) (-0.129,0.1073) (0.014,0.1219) (0.157,0.1366) (0.300,0.1510) (0.443,0.1643) (0.586,0.1757) (0.729,0.1845) (0.872,0.1901) (1.015,0.1920) (1.158,0.1901) (1.301,0.1845) (1.444,0.1757) (1.587,0.1643) (1.730,0.1510) (1.873,0.1366) (2.016,0.1219) (2.159,0.1073) (2.302,0.0935) (2.445,0.0806) (2.588,0.0690) (2.732,0.0586) (2.875,0.0494) (3.018,0.0415) (3.161,0.0347) (3.304,0.0289) (3.447,0.0240)
};
\fill[black!12,fill opacity=0.55] (1.573,0) -- plot[smooth] coordinates {
(1.573,0.0360) (1.687,0.0434) (1.802,0.0520) (1.916,0.0623) (2.031,0.0741) (2.145,0.0878) (2.260,0.1034) (2.374,0.1210) (2.488,0.1402) (2.603,0.1610) (2.717,0.1828) (2.832,0.2049) (2.946,0.2265) (3.060,0.2464) (3.175,0.2636) (3.289,0.2768) (3.404,0.2851) (3.518,0.2880) (3.633,0.2851) (3.747,0.2768) (3.861,0.2636) (3.976,0.2464) (4.090,0.2265) (4.205,0.2049) (4.319,0.1828) (4.434,0.1610) (4.548,0.1402) (4.662,0.1210) (4.777,0.1034) (4.891,0.0878) (5.006,0.0741) (5.120,0.0623) (5.235,0.0520) (5.349,0.0434) (5.463,0.0360)
} -- (5.463,0) -- cycle;
\draw[densely dashed,semithick] plot[smooth] coordinates {
(1.573,0.0360) (1.687,0.0434) (1.802,0.0520) (1.916,0.0623) (2.031,0.0741) (2.145,0.0878) (2.260,0.1034) (2.374,0.1210) (2.488,0.1402) (2.603,0.1610) (2.717,0.1828) (2.832,0.2049) (2.946,0.2265) (3.060,0.2464) (3.175,0.2636) (3.289,0.2768) (3.404,0.2851) (3.518,0.2880) (3.633,0.2851) (3.747,0.2768) (3.861,0.2636) (3.976,0.2464) (4.090,0.2265) (4.205,0.2049) (4.319,0.1828) (4.434,0.1610) (4.548,0.1402) (4.662,0.1210) (4.777,0.1034) (4.891,0.0878) (5.006,0.0741) (5.120,0.0623) (5.235,0.0520) (5.349,0.0434) (5.463,0.0360)
};
\fill[black!22,fill opacity=0.55] (2.531,0) -- plot[smooth] coordinates {
(2.531,0.0857) (2.631,0.1032) (2.732,0.1239) (2.832,0.1482) (2.932,0.1765) (3.032,0.2092) (3.132,0.2463) (3.232,0.2880) (3.332,0.3339) (3.432,0.3834) (3.532,0.4353) (3.633,0.4879) (3.733,0.5393) (3.833,0.5867) (3.933,0.6275) (4.033,0.6590) (4.133,0.6789) (4.233,0.6857) (4.333,0.6789) (4.434,0.6590) (4.534,0.6275) (4.634,0.5867) (4.734,0.5393) (4.834,0.4879) (4.934,0.4353) (5.034,0.3834) (5.134,0.3339) (5.235,0.2880) (5.335,0.2463) (5.435,0.2092) (5.535,0.1765) (5.635,0.1482) (5.735,0.1239) (5.835,0.1032) (5.935,0.0857)
} -- (5.935,0) -- cycle;
\draw[semithick] plot[smooth] coordinates {
(2.531,0.0857) (2.631,0.1032) (2.732,0.1239) (2.832,0.1482) (2.932,0.1765) (3.032,0.2092) (3.132,0.2463) (3.232,0.2880) (3.332,0.3339) (3.432,0.3834) (3.532,0.4353) (3.633,0.4879) (3.733,0.5393) (3.833,0.5867) (3.933,0.6275) (4.033,0.6590) (4.133,0.6789) (4.233,0.6857) (4.333,0.6789) (4.434,0.6590) (4.534,0.6275) (4.634,0.5867) (4.734,0.5393) (4.834,0.4879) (4.934,0.4353) (5.034,0.3834) (5.134,0.3339) (5.235,0.2880) (5.335,0.2463) (5.435,0.2092) (5.535,0.1765) (5.635,0.1482) (5.735,0.1239) (5.835,0.1032) (5.935,0.0857)
};
\fill[black!30,fill opacity=0.55] (4.026,0) -- plot[smooth] coordinates {
(4.026,0.2000) (4.112,0.2408) (4.198,0.2891) (4.283,0.3458) (4.369,0.4119) (4.455,0.4880) (4.541,0.5747) (4.627,0.6720) (4.712,0.7791) (4.798,0.8945) (4.884,1.0156) (4.970,1.1385) (5.056,1.2583) (5.142,1.3690) (5.227,1.4642) (5.313,1.5377) (5.399,1.5841) (5.485,1.6000) (5.571,1.5841) (5.656,1.5377) (5.742,1.4642) (5.828,1.3690) (5.914,1.2583) (6.000,1.1385) (6.086,1.0156) (6.171,0.8945) (6.257,0.7791) (6.343,0.6720) (6.429,0.5747) (6.515,0.4880) (6.600,0.4119) (6.686,0.3458) (6.772,0.2891) (6.858,0.2408) (6.944,0.2000)
} -- (6.944,0) -- cycle;
\draw[thick] plot[smooth] coordinates {
(4.026,0.2000) (4.112,0.2408) (4.198,0.2891) (4.283,0.3458) (4.369,0.4119) (4.455,0.4880) (4.541,0.5747) (4.627,0.6720) (4.712,0.7791) (4.798,0.8945) (4.884,1.0156) (4.970,1.1385) (5.056,1.2583) (5.142,1.3690) (5.227,1.4642) (5.313,1.5377) (5.399,1.5841) (5.485,1.6000) (5.571,1.5841) (5.656,1.5377) (5.742,1.4642) (5.828,1.3690) (5.914,1.2583) (6.000,1.1385) (6.086,1.0156) (6.171,0.8945) (6.257,0.7791) (6.343,0.6720) (6.429,0.5747) (6.515,0.4880) (6.600,0.4119) (6.686,0.3458) (6.772,0.2891) (6.858,0.2408) (6.944,0.2000)
};
\node[font=\scriptsize,align=left] at (-1.0,1.35) {$\dd Q/\dd V$\\[-1pt]{\tiny(arb.\ u.,}\\[-2pt]{\tiny$Q_j/4w_j$ scale)}};
\draw[-{Latex[length=1.4mm]}] (7.0,-0.85) -- (-1.4,-0.85) node[pos=0.62,below,font=\tiny] {delithiation (discharge): main direction, $\theta_j:1\to0$};
\end{scope}
\end{tikzpicture}
\caption{흑연 staging 의 갤러리 채움 도식(위)과 그에 대응하는 진짜-축척 전위 축(아래). 위: 수평선이
그래핀 층, 음영 띠가 리튬이 든 층간 갤러리다(stage $2\mathrm L$ 은 면내 질서가 풀린 단계라 옅게 표시).
채울수록 stage 번호가 줄어 stage 1($\mathrm{LiC_6}$)에 닿는다. 아래: 네 전이의 초기값 $U$(표~\ref{tab:staging},
0.210/0.140/0.120/0.085 V)를 \emph{실제} 전위 간격대로 배치하고, 각 위치에 그 전이가 남기는 $\dd Q/\dd V$
peak(식~\eqref{eq:belliden} 종, 높이 $Q_j/4w_j$, 폭 $w_j$ — 좌표는 T2\_calc.py 의 식 수치 평가)를 얹었다.
점선이 위 도식의 균등 칸 배치가 아래의 불균등 실제 전위 간격과 다름을 보인다 — $3\!\to\!2\mathrm L$ 과
$2\mathrm L\!\to\!2$ 는 실제로 $20$ mV 밖에 떨어져 있지 않아 peak 가 서로 겹친다(음영 중첩 $=$ 식~\eqref{eq:sum}
합산의 시각화). peak 높이는 stage 1 방향으로 커진다($Q_j$ 증가$\cdot$$w_j$ 감소).}
\label{fig:staging}
